\begin{pastexam}[label={px:ntu-110-ie-ds-02}]{臺大工工所 110 學年度 第 2 題}
（單選；每題 3 分，答錯不倒扣）

In this exam, the height of a leaf (in either a tree or heap) is defined as 0.

For an AVL tree with 2021 nodes, let x be the minimum possible height, y be the maximum height, then
(y-x) mod 5 is

(A)~0 \qquad (B)~1 \qquad (C)~2 \qquad (D)~3 \qquad (E)~4
\end{pastexam}

\begin{pastsoln}
答案是 (E)。 $x = 10$、$y = 14$，$(y - x) \bmod 5 = 4$。

本卷的高度以邊數計（葉節點高度 $0$），比資結書數節點的高度少 $1$。

\textbf{最小高度 $x$}。 高度 $h$ 的二元樹最多 $2^{h+1} - 1$ 個節點。 $2^{10} - 1 = 1023 < 2021 \le 2047 = 2^{11} - 1$，
所以高度至少 $10$。 高度 $10$ 做得到：把 $2021$ 個鍵依序排好，取中間當根、左右遞迴，建出的完全平衡樹
高度 $\lfloor \lg 2021 \rfloor = 10$，每個節點左右高度差不超過 $1$，是 AVL 樹。 $x = 10$。

\textbf{最大高度 $y$}。 令 $N(h)$ 是高度 $h$ 的 AVL 樹最少的節點數。 最瘦的樹是根加上高度 $h - 1$ 和 $h - 2$ 的
最瘦子樹（費氏樹，資結書第~10.2.2~節），所以
\[
  N(0) = 1,\quad N(1) = 2,\quad N(h) = N(h-1) + N(h-2) + 1 .
\]
（換成資結書的記號就是 $N(h) = F_{h+3} - 1$。） 依序算出：
\begin{center}
\begin{tabular}{c*{8}{c}}
\toprule
$h$ & 8 & 9 & 10 & 11 & 12 & 13 & 14 & 15 \\
$N(h)$ & 88 & 143 & 232 & 376 & 609 & 986 & 1596 & 2583 \\
\bottomrule
\end{tabular}
\end{center}
（$h = 0..7$ 是 $1, 2, 4, 7, 12, 20, 33, 54$。） 高度 $14$ 最少要 $1596 \le 2021$ 個節點，而且做得到：高度 $14$
的 AVL 樹，節點數可以是 $N(14) = 1596$ 到 $2^{15} - 1$ 之間的任何整數。 （對高度歸納：高度 $h$ 的 AVL 樹，
左右子樹的高度是 $h - 1, h - 1$ 或 $h - 1, h - 2$；由歸納假設各自的節點數可以取遍一個區間，加起來的範圍從
$N(h - 1) + N(h - 2) + 1 = N(h)$ 一路連到 $2(2^h - 1) + 1 = 2^{h+1} - 1$，中間沒有空隙。） 所以有 $2021$ 個
節點、高度 $14$ 的 AVL 樹。 高度 $15$ 至少要 $2583 > 2021$ 個節點，做不到。 所以 $y = 14$。

$(y - x) \bmod 5 = 4 \bmod 5 = 4$，選 (E)。 如果用數節點的高度（葉節點高度 $1$），$x$、$y$ 都加 $1$，差不變，
答案一樣。
\end{pastsoln}
